Q3与方法78逐局相同，因此其对方法56的1.48秒收益全部来自既有的安全提前停车，不能记为候选84带来的新增规划收益。Q4相对方法78平均节省28.41秒、0.496\%，95\%区间为[-5.92,62.75]秒，86局变慢；相对方法56则为29.62秒、0.516\%，区间[-4.69,63.92]秒，81局变慢。两个对照组（56与78）在Q4的区间均跨零，按预先声明的推广规则不推广候选84，默认仍保留方法78；表中凡涉及区间的行均已在(a)(b)两栏分别标明对照组。最大退步出现在同一个16源环境（种子2610038），候选84比方法78多耗708.96秒、比方法56多耗707.85秒；最大收益也出现在同一环境下两法相同（种子2610090，1635.82秒）。这说明Q4的收益与退步都由少数高源数环境主导，均值不足以代表逐局表现。实验候选可由geometric策略名显式运行。
